\chapter{Confronto fra configurazioni, e l'angolo di ingresso}
\label{ch:compare}

Tutto fin qui ha certificato una macchina alla volta. Questo capitolo
pone la domanda che un progettista si fa davvero --- campana o plug,
quanto lungo, quanto swirl, a quale angolo di ingresso --- e scopre
che la parte difficile non è calcolare le risposte ma renderle
\emph{confrontabili}. Registra anche la più lunga svolta sbagliata di
questo programma: quattro sessioni spese a inseguire un difetto che
viveva in uno strumento di misura anziché nel flusso.

\section{La domanda, e perché è soprattutto una questione di
contabilità}

La spinta scala con la portata. Una campana calcolata in un mondo
posto in un certo modo, messa contro un plug calcolato in un altro,
confronta i \emph{mondi}, non i tipi di ugello. Perciò, prima che
qualunque confronto abbia senso, va fissata e imposta una base
comune:

\begin{itemize}[itemsep=1pt]
\item lo stesso gas, lo stesso stato di camera $(P_0, T_0)$, lo
  stesso ambiente $p_a$;
\item la stessa \textbf{portata} --- che porta un bonus, perché per
  una gola bloccata $\dot m = p_c A_t/c^{*}$, quindi portata uguale a
  stato di camera uguale significa \emph{area di gola} uguale, e
  perciò il coefficiente di spinta $C_F = F/(p_c A_t)$ e l'impulso
  specifico sono numeri direttamente confrontabili con la dimensione
  della macchina divisa via;
\item lo stesso \textbf{ingombro}: a ogni configurazione è concesso
  lo stesso raggio massimo. Un plug cui si conceda raggio illimitato
  vince banalmente sulla sola espansione.
\end{itemize}

Ogni configurazione è poi ottimizzata rispetto alla \emph{propria}
variabile di progetto: la campana rispetto al rapporto di espansione,
i plug rispetto al punto in cui la spina viene tagliata.

Una sottigliezza rende i due lati la stessa grandezza fisica. Per un
razzo alimentato da una camera chiusa la spinta assiale eguaglia il
flusso di quantità di moto in gauge ambiente attraverso \emph{una
qualunque} superficie che copra l'intero scarico e racchiuda
l'hardware. Per la campana quella superficie è il piano di uscita.
Per il plug lo stesso teorema si applica alla linea iniziale della
marcia: il flusso lì contiene già la spinta di tutto ciò che sta a
monte, perché la frontiera di getto libero fra il labbro e quel
taglio contribuisce esattamente zero nel gauge ($p = p_a$ su di
essa). Quindi
$F_{\rm plug}(l) = F_{\rm in}(x_0) + \int_{x_0}^{l}(p_w - p_a)\,
2\pi y\,(-\mathrm{d}y)$ è la spinta totale dalla camera fino al
troncamento --- lo stesso oggetto di quella della campana.

\section{Risultati di riferimento}

\begin{table}[htbp]
\centering
\small
\begin{tabular}{lrrrr}
\toprule
configurazione & $F$ [MN] & $C_F$ [--] & $I_{sp}$ [s] & $L$ [m] \\
\midrule
campana, al suo ottimo di rapporto & 116.474 & 1.4654 & 279.62 & 7.34 \\
plug, spina completa               & 116.148 & 1.4613 & 278.84 & 6.00 \\
plug, troncato                     & 116.148 & 1.4613 & 278.84 & 3.69 \\
plug, troncato $+$ swirl           & 111.839 & 1.4071 & 268.49 & 2.56 \\
\bottomrule
\end{tabular}
\caption{Quattro configurazioni a un punto di funzionamento sulla
base comune della Sezione~1: stesso gas, camera, ambiente, portata
(e quindi area di gola) e raggio di ingombro. Spinta in meganewton,
coefficiente di spinta e impulso specifico adimensionale e in
secondi, lunghezza in metri.}
\label{tab:configcompare}
\end{table}

\textbf{La campana vince dello $0.282\%$, e il risultato è
risolto.} La banda derivata sulla spinta viene dal raffinamento della
marcia stessa: $J = 1.161610, 1.161481, 1.161454 \times 10^{8}$~N a
tre risoluzioni, differenze successive $1.287\times10^{4}$ e poi
$2.737\times10^{3}$ --- un rapporto di $4.70$, che conferma il
secondo ordine --- da cui una banda di Richardson di
$1.09\times10^{4}$~N, ovvero lo $0.0094\%$. Il divario
campana--plug di $3.28\times10^{5}$~N vale $30\times$ quella banda.

\textbf{Il plug eguaglia la campana in metà lunghezza}, 3.69~m contro
7.34~m, che è il risultato dell'aerospike enunciato nel modo in cui
la termodinamica permette di enunciarlo
(Sezione~\ref{sec:ceiling}).

\textbf{Lo swirl costa il $4.0\%$}, un ordine di grandezza più di
qualunque differenza geometrica della tabella --- ma la spina con
swirl scende fino a $38\times$ sotto l'ambiente, un regime in cui un
flusso reale separerebbe e questa marcia non viscosa non può
seguirlo. Quel numero è indicativo, ed è etichettato come tale.

\begin{figure}[htbp]
\centering
\includegraphics[width=\textwidth]{fig_config_compare.pdf}
\caption{Le configurazioni a confronto. (a)~L'hardware in un'unica
scala, posizione assiale e raggio in metri: la parete della campana
al proprio ottimo, la spina del plug adattata in massa, il suo
troncamento e la spina ridimensionata per lo swirl. (b)~Pressione di
parete in bar contro la linea dell'ambiente: sopra di essa la parete
fa spinta, nella regione ombreggiata sotto la parete tira
all'indietro --- e la spina con swirl entra molto dentro quella
regione. (c)~Spinta in meganewton contro lunghezza dell'ugello: la
curva del plug sale e satura, raggiungendo la spinta della campana
(blu) a circa metà della sua lunghezza. La banda ombreggiata è la
banda di Richardson \emph{derivata} sulla spinta, $\pm0.009\%$; il
divario campana--plug vale $30\times$ essa.}
\label{fig:configcompare}
\end{figure}

\section{Perché sono così vicine: il tetto termodinamico}
\label{sec:ceiling}

La quasi uguaglianza dei coefficienti non è una coincidenza né una
delusione. Entrambi gli ugelli espandono lo stesso gas dalla stessa
camera alla stessa pressione di uscita, e la sola conservazione
dell'energia fissa allora la velocità di uscita:
$\tfrac12 q_e^2 = h_0 - h(p_a)$ dà $q_e = 2742.136$~m/s senza alcuna
geometria nell'enunciato. Poiché la base impone la stessa portata, la
spinta ideale di quantità di moto
$\dot m\,q_e = 1.16474\times10^{8}$~N è esattamente la spinta
misurata della campana --- un'identità, perché al proprio ottimo la
campana è perfettamente espansa (il termine di pressione si annulla)
e la sua uscita è esattamente assiale (misurata a
$\theta_e = -10^{-16}$ nel Capitolo~\ref{ch:score}).

La campana sta \emph{sul} tetto. La geometria può quindi soltanto
perdere rispetto a essa, per divergenza, non uniformità ed espansione
incompleta --- tutti effetti dell'ordine dell'uno percento. È tutto
qui il campo da gioco dell'ottimizzazione di contorno, ed è il motivo
per cui questo programma deriva ogni tolleranza invece di
sceglierne una.

\section{Un vuoto di certificazione che il confronto ha esposto}

Predisporre il lato plug ha richiesto di porre il flusso come fa un
vero aerospike: arrivando al labbro puntato \emph{verso l'interno}
così che il ventaglio lo riporti verso l'assiale, con la spina
discendente verso l'asse. Ogni mondo di plug assialsimmetrico
certificato in questo documento pone l'opposto --- ingresso assiale,
il ventaglio che lo ruota \emph{verso l'esterno}, la ``spina'' che si
alza allontanandosi dall'asse (Figura~\ref{fig:turndir}a).

Non sono immagini speculari. Il termine sorgente assialsimmetrico è
$\delta\,c^2 v/y$ e $y$ è sempre positivo, quindi
$\theta \to -\theta$ non è una simmetria: il caso verso l'esterno
corre con $v > 0$ e $y$ crescente, quindi la sorgente si indebolisce
lungo la marcia, mentre il caso verso l'interno corre con $v < 0$ e
$y$ decrescente, quindi si rafforza. L'oracolo piano è immune
(nessuna sorgente, simmetria speculare genuina), quindi è
specificamente la certificazione \emph{assialsimmetrica} ad aver
esercitato sempre e solo un lato.

La rotazione verso l'esterno non è mai stata un progetto candidato
--- era un banco di prova. Un plug che ruota verso l'esterno lascia
l'intero scarico a $\theta_E > 0$ e raccoglie solo $\cos\theta_E$
della quantità di moto: $10.6\%$ perso alla rotazione di progetto di
$26.7^{\circ}$, contro nulla per un progetto verso l'interno. È il
modo naturale di scrivere un campo di prova a getto di spigolo in
forma chiusa, ed è così che è diventato l'unico caso testato.

\begin{figure}[htbp]
\centering
\includegraphics[width=\textwidth]{fig_turn_direction.pdf}
\caption{Rotazione verso l'esterno contro rotazione verso l'interno,
posizione assiale e raggio in metri. (a)~Ciò che ogni suite
assialsimmetrica certificata poneva: il flusso arriva assiale, il
ventaglio lo ruota verso l'esterno e la ``spina'' si alza
allontanandosi dall'asse; $v>0$, $y$ cresce, la sorgente si
indebolisce. (b)~Ciò che un vero aerospike è, e ciò che questo
capitolo pone: il flusso arriva puntato verso l'asse, il ventaglio lo
riporta indietro, la spina scende; $v<0$, $y$ decresce, la sorgente
si rafforza. (c)~La conseguenza misurata, in termini di non chiusura
di massa della marcia contro l'angolo di rotazione del ventaglio: la
rotazione verso l'esterno resta al pavimento senza sorgente
(punteggiato) e rimane piatta, quella verso l'interno sale a
$2\times10^{-2}$.}
\label{fig:turndir}
\end{figure}

\section{L'oracolo del flusso radiale}
\label{sec:radialoracle}

Per scoprire se quella differenza fosse un difetto, il programma
aveva bisogno di una soluzione assialsimmetrica esatta con il flusso
diretto \emph{verso} l'asse. Il flusso radiale sferico ne fornisce
una: ogni linea di corrente è un raggio uscente da un punto
sull'asse, la velocità dipende solo dalla distanza da esso, e la
continuità con l'isentropica la chiude come
$\rho(q)\,q\,r^2 = \mathrm{cost}$. Mettendo il punto singolare a
monte si ottiene un flusso divergente ($v>0$); mettendolo a valle uno
convergente ($v<0$). Un'unica formula, quindi i due casi non
differiscono in nulla se non nel segno che le equazioni vedono. E
poiché le linee di corrente sono raggi, un cono è una parete
\emph{esatta}.

\carrier{validation/a1\_source\_flow\_oracle.py}, \textbf{6/6 PASS.}
Il bordo libero è escluso per costruzione: entrano nel confronto solo
i punti sotto un limite rigoroso di dipendenza --- una retta dalla
sommità della linea iniziale con la pendenza caratteristica più
ripida presente nel campo, che nulla proveniente dall'alto può
superare. Rispetto alla forma chiusa, gli errori mediani sono

\begin{center}
\small
\begin{tabular}{lrr}
\toprule
 & convergente ($v<0$) & divergente ($v>0$) \\
\midrule
pressione di parete & $7.9\times10^{-4}$ & $5.2\times10^{-4}$ \\
interno $|\Delta q|/q$ & $3.4\times10^{-4}$ & $1.8\times10^{-4}$ \\
\bottomrule
\end{tabular}
\end{center}

\noindent Le celle interne e quelle a parete assegnata sono pulite
per flusso che ruota verso l'interno. Una trappola utile è emersa
costruendo il controllo negativo: nel flusso radiale \emph{ogni} cono
è una linea di corrente esatta, quindi perturbare l'angolo del cono
produce semplicemente un'altra soluzione esatta e non può mai
rigettare. Il rigettatore valido incurva la parete allontanandola dal
raggio, rompendo la proprietà di raggio, e separa di $29\times$.
Anche i rigettatori vanno scelti in base alla fisica --- la stessa
lezione dell'Appendice~\ref{app:f2rejector}, incontrata di nuovo.

\section{Lo strumento, non il flusso}

La non chiusura di massa che ha dato origine a questa indagine è in
larga parte una proprietà del misuratore. La verifica di
conservazione integrava sull'\emph{ultima colonna} della marcia, che è
una caratteristica, non un piano, e le cui lunghezze di segmento sono
selvaggiamente disuguali: su una spina discendente la più lunga è
$0.84$~m contro una media di $0.064$~m. L'errore trapezoidale scala
con il quadrato del segmento, quindi quell'unico segmento dominava.
Tre misure risolvono la questione: il flusso proprio della linea di
corrente del bordo libero è $4\times10^{-11}$ --- zero di macchina,
quindi non perde nulla; la parete perde lo $0.016\%$; e le celle
riproducono la soluzione esatta di cui sopra. L'ispezione colonna per
colonna trova la causa: solo la \emph{prima} colonna marciata ha un
buco di copertura, $26\times$ la spaziatura media fra righe quando la
parete sale e $39\times$ quando scende, contro $0.5$--$1.0\times$ per
ogni altra colonna --- perché una colonna è una caratteristica mentre
una linea iniziale verticale non lo è, quindi la prima colonna deve
da sola fare da ponte fra le due geometrie.

La lezione è quella che vale la pena portarsi dietro: \textbf{quando
la grandezza di interesse ha una banda derivata disponibile in modo
diretto, la si derivi in modo diretto.} L'incertezza sulla spinta era
stata dedotta da un misuratore di conservazione che misurava
qualcos'altro; raffinare la marcia e lasciarle misurare il proprio
errore --- la prima regola stessa del programma --- ha risolto la
questione in una sola esecuzione.

\section{L'angolo di ingresso come variabile di progetto}
\label{sec:inletangle}

Per direttiva del proprietario l'angolo di rotazione iniziale non è
più posto: fa parte dell'ottimizzazione. L'RDE consegna il proprio
flusso attraverso un ingresso anulare \emph{rettilineo} fra il raggio
iniziale della spina e il labbro del cowl, e l'inclinazione di quel
ingresso è la variabile di progetto.

La rotazione del ventaglio di labbro $\Delta\nu$ è fissata
dall'espansione (da $M_i$ fino a $p_a$) e non dipende dall'angolo di
ingresso, quindi l'angolo di scarico è semplicemente
$\theta_E = \theta_i + \Delta\nu$; la spazzata è riportata in
$\theta_E$ perché è quello che interessa alla fisica, e la classica
condizione di progetto è $\theta_E = 0$. La base della Sezione~1 è
invariata, con un'aggiunta: il flusso assiale di massa attraverso un
ingresso inclinato scala con $\cos\theta_i$, quindi l'anello viene
ridimensionato a ogni candidato per mantenere la portata di
riferimento, e la spinta è confrontata a \emph{lunghezza fissa} così
che la lunghezza non venga barattata in silenzio contro l'angolo.

\carrier{validation/a1\_inlet\_angle\_opt.py}. La spazzata alla cieca
atterra a $\theta_E = 0$ --- la previsione della divergenza,
recuperata anziché imposta --- e la penalità misurata segue il
modello $\cos\theta_E$ in segno e ordine su tutta la finestra
ammissibile:

\begin{center}
\small
\begin{tabular}{rrr}
\toprule
$\theta_E$ [gradi] & misurata & modello $\cos\theta_E$ \\
\midrule
$-6$ & $-0.365\%$ & $-0.548\%$ \\
$-4$ & $-0.153\%$ & $-0.244\%$ \\
$-2$ & $-0.031\%$ & $-0.061\%$ \\
$0$  & migliore   & --- \\
$+2$ & $-0.068\%$ & $-0.061\%$ \\
$+4$ & $-0.224\%$ & $-0.244\%$ \\
$+8$ & $-0.799\%$ & $-0.973\%$ \\
$+12$& $-1.753\%$ & $-2.185\%$ \\
\bottomrule
\end{tabular}
\end{center}

Due cose vanno lette. La penalità misurata è sistematicamente
\emph{minore} della pura divergenza, perché la distribuzione di
pressione di parete compensa in parte --- il baratto che il modello
in forma chiusa non contiene. E la curva è asimmetrica: ruotare verso
l'interno oltre l'ottimo costa meno che ruotare verso l'esterno dello
stesso angolo, che è la forma quantitativa della premessa della
direttiva secondo cui gli aerospike che ruotano verso l'esterno non
sono competitivi. La finestra è limitata dal lato interno dalla
geometria, non da una preferenza: a $\theta_E = -8^{\circ}$ la spina
scende a $y = 0.32$~m e la marcia perde la certificazione mentre il
flusso corre verso l'asse, dove la sorgente assialsimmetrica diverge.
